\section{Limitations}\label{sec:limitations}

\textbf{Evaluation scope.} Models are evaluated with checkers disabled, and defenses only on SD1.4, since weight editing needs architecture-specific DiT ports.

\noindent
\textbf{Coverage.} \looseness=-1 Construction models and closed commercial systems are excluded. Severity definitions are Western-centric (a non-Western probe is in \SuppSec{\ref{app:cultural}}), and prompts are English-only. Copyright and PID show lower HGR than NSFW (\Tab{\ref{tab:t2i_models}}) because MLLM classifiers are weaker at identity harms, so the gap there is likely \emph{underestimated}. NSFW-targeted filters and single-category detectors (e.g., NudeNet~\citep{chhabra2020nudenet}) are out-of-domain on Legal and Social content (Insight~5). The human study has three annotators; a larger panel and a fully blind harm rating are planned.

\noindent
\textbf{Evaluator error.} \looseness=-1 The evaluator is noisiest at L1 (human $\kappa = 0.73$; guidance-ablation shift $<2.4\%$). Because it sees sample-specific risk descriptions, absolute HGR figures are protocol-relative (at L3, 7.2\% blind vs.\ 76.4\% guided), and L0 rates are the noise floor. The label-stripping ablation retains the sample-specific descriptions, so description-cue bias is not fully excluded; an image-grounded audit is future work.
The ladder scorer and verifier are Qwen-family, but the InternVL2-8B run (Insight~6) shows the transfer is not a same-family artifact. Correlation within a ladder also understates the reported SEMs.
